\documentclass[12pt,a4paper]{book}

% Formato del documento
%\usepackage[papersize={210mm,297mm},inner=3.5cm,outer=2cm,top=2.5cm,bottom=2.5cm]{geometry}
%\renewcommand{\baselinestretch}{1}
%\setlength{\parskip}{8pt}

\usepackage[utf8]{inputenc}
\usepackage{amssymb}
\usepackage{amsmath}
\usepackage{graphicx}
\usepackage{tikz}
\usepackage{tkz-graph}
\usepackage{hyperref}
\usepackage[parfill]{parskip}
\usepackage{float}

\renewcommand{\chaptername}{Capítulo}
\renewcommand{\contentsname}{Índice}
\renewcommand{\bibname}{Bibliografía}
\renewcommand{\figurename}{Figura}

\definecolor{verde_oscuro}{rgb}{0.0, 0.5, 0.0}

\newtheorem{defi}{Definición}[section]
\newtheorem{prop}{Proposición}[section]
\newtheorem{propi}{Propiedades}[section]
\newtheorem{lema}{Lema}[section]
\newtheorem{tma}{Teorema}[section]
\newtheorem{cor}{Corolario}[section]
\newtheorem{nota}{Nota}[section]
\newtheorem{ejem}{Ejemplo}[section]

\hypersetup{
    colorlinks=true,
    linkcolor=blue,
    filecolor=magenta,      
    urlcolor=cyan,
    pdftitle={Overleaf Example},
    pdfpagemode=FullScreen,
    }

\urlstyle{same}

%%%%%%%ALTERNATIVA 2%%%%%%%%%%%%
\textheight=21cm
\textwidth=17cm
%\topmargin=-1cm
\oddsidemargin=0cm
\evensidemargin=0cm
\parindent=0mm
\pagestyle{plain}
%%%%%%%%%%%%%%%%%%%%%

%%%%%%%  Comando Portada  %%%%%%%%%%%
\newcommand{\nuevaportada}[6]{
    \thispagestyle{empty}
    \begin{center}
        \includegraphics[width=0.5\textwidth]{images/logo.png}
        
        \vspace{0.5cm}
        {\Large\bfseries\textsc{M\'aster Universitario en #1}\par}
        
        \vspace{0.5cm}
        \includegraphics[width=0.4\textwidth]{images/uv.png}
        
        \vspace{0.5cm}
        {\Large\bfseries\textsc{Trabajo de Fin de M\'aster}\par}
        
        \vspace{0.5cm}
        {\Large\bfseries #2\par}
        
        \vspace{2cm}
        \begin{flushright}
            \begin{tabular}{l} 
                {\large\bfseries\textsc{Autor:}} \\
                {\large\textsc{#3}} \\ [0.5cm] 
                {\large\bfseries\textsc{Tutora:}} \\ 
                {\large\textsc{#4}} \\ [0.5cm]
                {\large\bfseries #5} 
            \end{tabular}
        \end{flushright}
    \end{center}
    \clearpage
}

\begin{document}

\nuevaportada{Ciencia de Datos}{Problema de Localización de Centros k-Balanceado Multiobjetivo}{Manuel Rubio Martínez}{Anna Martínez Gavara}{Abril, 2025}

\clearpage

\newpage
\tableofcontents

\chapter{Preliminares}
\section{primera parte}

\begin{defi}
Un algoritmo \textbf{"greedy"} (codicioso) es un método construcción de soluciones para un problema de optimización,
de manera que en cada paso se escoge la mejor posible, aunque esto no lleve finalmente a un óptimo global.
\end{defi}

\bigskip

\begin{defi}
Un algoritmo \textbf{"greedy randomizado"} se basa en el "greedy", pero en cada paso se elige de forma aleatoria entre un conjunto de mejores opciones.
Si el conjunto de opciones es el total, el algoritmo es "random".
\end{defi}

\chapter{Introducción}

\section{Antecedentes}
\subsection{Optimización Combinatoria}

Es la rama de la optimización enfocada en encontrar óptimos en conjuntos de soluciones finitos.
Esto no significa que el problema sea fácil, al contrario, por lo general es un problema NP-complejo. \\
Que sea NP-complejo implica que conforme el número de soluciones aumenta, la complejidad computacional lo hace, a su vez, exponencialmente. \\
Un ejemplo de ello es el problema sobre el que trata este TFM, 
tratar de que encontrar, en un conjunto de $n$ elementos, un subconjunto $k$ de ellos tal que  minimice varias funciones objetivo.
\\ \\
Es rápido calcular las combinaciones posibles (siempre suponiendo $k \leq  n$) son:

$$
    \binom{n}{k} = \frac{n!}{(n-k)!k!}
$$

Por ejemplo, si $n=50$ y $k=5$, entonces
$$
    \frac{50!}{(50-5)!5!}=2118760
$$
para un problema bastante sencillo. Para casos reales, podríamos tener n=5000 y k=50, la cantidad de posibles solucones sería del orden de $2.28  e^{120}$. Contabilizando el coste de 
evaluación de cada solución, no es viable encontrar el óptimo de forma exacta. Incluso utilizando los ordenadores más potentes, estaríamos hablando de tiempos astronómicos.

\subsection{Metaheurísticos}

Las metaheurísticas aparecen con el desarrollo de las ciencias de la computación, como principal solución a estos problemas. Donde los métodos exactos no serían capaces de encontrar el óptimo,
los Metaheurísticos buscan soluciones `buenas` en tiempos de ejecución razonables, sacrificando la garantía que proporcionan los métodos exactos de obtener un óptimo global.

Generalmente combinan heurísticas para obtener soluciones buenas, con aleatoriedad, para conseguir una exploración
más ámplia del conjunto de soluciones posibles, buscando un balance entre exploración y explotación del entorno de soluciones.\\
Exploración ya que no se quedan bloqueados en regiones aisladas, intentando buscar óptimos locales, sino que recorren el resto del espacio buscando nuevas soluciones.\\
Explotación porque cuando entran en una región, profundizan en ella buscando los mejores óptimos locales.

Son algoritmos muy generales, fáciles de implementar y adaptables a casi 
cualquier problema de optimización.

Los principales metaheurísticas son:
\begin{itemize}
    \item Algoritmos Genéticos (GA - Genetic Algorithms)\\
    Inspirados en la selección natural y el funcionamiento de la genética, se basan en la generación aleatoria de soluciones, y el cruce de las más prometedoras entre sí, 
    añadiendo mutaciones. El cruce de las más prometedoras crea soluciones buenas, mientras las mutaciones aumentan la diversidad.
    \item Recocido Simulado (SA - Simulated Annealing)\\
    Basada en la física, se basa en el proceso de recocer metales (calentarlos y enfriarlos para alterar su estructura). Explora soluciones aceptando movimientos de mejora, pero también que puedan empeorar, 
    en función de una probabilidad dependiente de la "temperatura", que decrece con el paso de las iteraciones. Al principio, una temperatura alta favorece la exploración, y al final una baja favorece la explotación.
    \item Búsqueda Tabú (TS - Tabu Search)
    Genera una lista `Tabú` que actúa como memoria del algoritmo, incluyendo los movimientos recientes con el objetivo de no repetir soluciones recientes, sino explorar nuevas, aunque estas no sean las mejores posibles. Esto favorece la exploración.
    \item Optimización por Colonia de Hormigas (ACO - Ant Colony Optimization)
    Diseñado para problemas de rutas, construyendo soluciones iterativamente a través de un grafo, cuyas aristas  se van pesando a medida que generan soluciones buenas.
    %\item Optimización por Enjambre de Partículas (PSO - Particle Swarm Optimization)
\end{itemize}


\subsection{GRASP (Greedy Randomized Adaptive Search Procedure)}

\section{Problemas de localización}
Aquí es clave. Describe los problemas típicos de localización:

Los problemas de localización buscan determinar la mejor ubicación para instalar facilidades (almacenes, centros de distribución, hospitales, etc.) para optimizar criterios como distancia a los clientes, coste total o cobertura. Ejemplos clásicos incluyen:

Problema del p-mediana: minimizar la distancia total entre clientes y facilidades.

Problema del p-center: minimizar la máxima distancia entre un cliente y su centro asignado.

Problema de localización de instalaciones no capacitado (Uncapacitated Facility Location Problem, UFLP): elegir qué facilidades abrir y a qué clientes servir, minimizando costes de apertura y asignación.

Aquí puedes incluir una breve justificación de por qué estos problemas son relevantes en la práctica (logística, salud, servicios públicos…).

\subsection{Estado del arte}

Menciona estudios previos donde se ha aplicado GRASP a problemas de localización:

Numerosos trabajos han demostrado la eficacia de GRASP en problemas de localización. 
Por ejemplo, Feo y Resende (1995) aplicaron GRASP al problema del p-median con resultados competitivos. 
Desde entonces, se han desarrollado numerosas variantes híbridas que combinan GRASP con técnicas como Path Relinking o aprendizaje automático para mejorar su rendimiento.


\section{Motivación y caso práctico}

\subsection{Motivación}

En el contexto actual de distribución y venta minorista, las empresas que comercializan productos a través de comercios asociados se enfrentan a un problema logístico y estratégico cada vez más complejo: decidir en qué comercios ubicar la oferta de un determinado producto, de manera que se garantice una alta disponibilidad para los clientes, una repartición equitativa de la carga de trabajo entre los establecimientos, y se eviten situaciones de saturación de la demanda en puntos concretos.

La resolución de este problema implica equilibrar múltiples objetivos que pueden entrar en conflicto. Por ejemplo, concentrar las ventas en los comercios más céntricos puede facilitar el acceso para más clientes, pero puede llevar a la sobrecarga de dichos comercios. Por otro lado, una asignación equitativa del volumen de ventas podría perjudicar la proximidad al cliente final.

Este tipo de problemas, típicos en entornos logísticos y de gestión de redes comerciales, son ideales para abordar desde una perspectiva de optimización multiobjetivo. En particular, el uso de metaheurísticas como GRASP (Greedy Randomized Adaptive Search Procedure) permite generar soluciones de alta calidad en espacios de búsqueda complejos, combinando eficiencia, aleatoriedad y adaptabilidad. Esto resulta especialmente valioso en problemas reales donde las soluciones exactas no son computacionalmente viables.

\subsection{Caso práctico}

Se propone aplicar GRASP para resolver el siguiente problema práctico:

Una empresa desea lanzar un nuevo producto en una red de comercios asociados distribuidos geográficamente en una ciudad o región. La empresa no puede colocar el producto en todos los comercios, por lo que debe seleccionar un subconjunto óptimo de ellos. Esta selección debe cumplir los siguientes criterios:

Cobertura geográfica: Los comercios seleccionados deben estar lo más cerca posible de toda la población de clientes potenciales, minimizando la distancia media de los clientes al comercio más cercano.

Limitación del punto más cargado: Se desea minimizar la carga del comercio más exigido, es decir, evitar que uno solo concentre una proporción demasiado alta de las ventas.

Equidad en la carga de trabajo: Se busca que todos los comercios seleccionados reciban un volumen similar de clientes o ventas, para evitar que algunos estén infrautilizados mientras otros estén sobrecargados.

Este problema se formula como una optimización multiobjetivo discreta y se aborda mediante una implementación de GRASP adaptada a la estructura del problema. La solución final debería ofrecer al decisor un conjunto de configuraciones de comercios que representen distintos compromisos entre los objetivos, permitiendo una elección informada.

Este TFM contribuye tanto al avance práctico en la toma de decisiones comerciales como al desarrollo académico en el campo de la optimización heurística multiobjetivo.


\section{Otras aplicaciones}

1. Selección de ubicaciones para puntos de recogida de pedidos (e-commerce o logística inversa)
Objetivos:

Minimizar la distancia media desde los clientes.

Equilibrar el volumen de recogidas entre los puntos.

Minimizar la sobrecarga del punto con más actividad.

Uso de GRASP: Selección iterativa y adaptativa de ubicaciones con evaluaciones parciales del coste total.

2. Selección de ubicaciones para oficinas o delegaciones de una empresa
Objetivos:

Cubrir bien la demanda regional de servicios.

Minimizar costes de operación.

Repartir equitativamente la carga de trabajo regional.

3. Ubicación de estaciones de carga para vehículos eléctricos
Objetivos:

Maximizar cobertura geográfica.

Minimizar costes de instalación.

Equilibrar la demanda esperada entre estaciones.

Uso de GRASP: Selección iterativa de ubicaciones, con evaluación de cobertura y demanda prevista.

Uso de GRASP: Evaluación de diferentes combinaciones de oficinas candidatas y asignación óptima de áreas de influencia.
\chapter{Modelo Matemático}

\section{Definición}
Funciones objetivo, restricciones, etc

\chapter{Grasp}




\chapter{Experimentación y Resultados}

Instancias, comparativa, CASO DE ESTUDIO

\chapter{Conclusiones}





\begin{thebibliography}{X}
    \bibitem{Juegos} \textsc{Martín Novo, Eduardo; Méndez Alonso, Alfredo}.
    \textit{Aplicaciones de la teoría de grafos a algunos juegos de estrategia.}
    \url{https://redined.educacion.gob.es/xmlui/handle/11162/13874}

    \bibitem{Caníbales} Caníbales y misioneros: \url{https://culturacientifica.com/2016/05/04/problema-los-tres-caballeros-los-tres-criados/}

    \bibitem{Laberintos} Laberintos: \url{https://www.cantab.net/users/michael.behrend/repubs/maze_maths/pages/tarry_lab_en.html}
    
    \bibitem{Ajedrez} Sobre ajedrez: \url{https://chess24.com/es/informate/noticias/como-piensan-y-operan-los-modulos-de-analisis}

    \bibitem{Mantel} Teorema de Mantel: \url{https://lidicky.name/oldteaching/20.569X/l13%20-%20Mantel.pdf}

    \bibitem{Turan} Vascimini, Vincent. (2017). Extremal Graph Theory: Turán’s Theorem. 
    In BSU Honors Program Theses and Projects. Item 234. Available at: \url{https://vc.bridgew.edu/honors_proj/234}

    \bibitem{otros_teoremas} \textsc{Widdershoven, Cas.} \textit{Extremal graphs and the Erd{\H{o}}s-Stone-Simonovits theorems.}
    \url{https://studenttheses.uu.nl/handle/20.500.12932/23972}

\end{thebibliography}


\end{document}

